\documentclass[11pt]{article}

\usepackage{amsmath,amsthm,amssymb}
\usepackage[margin=1in]{geometry}
\usepackage{enumitem}
\usepackage{hyperref}
\hypersetup{colorlinks=true,linkcolor=blue,urlcolor=blue}

\theoremstyle{plain}
\newtheorem{theorem}{Theorem}
\newtheorem{proposition}{Proposition}
\newtheorem{corollary}{Corollary}

\theoremstyle{definition}
\newtheorem{definition}{Definition}[section]

\theoremstyle{remark}
\newtheorem{remark}{Remark}[section]

\newcommand{\Acat}{\mathcal{A}}
\renewcommand{\Im}{\operatorname{Im}}
\newcommand{\Ob}{\operatorname{Ob}}
\newcommand{\NK}{N_K}

\title{The closure criterion for handedness --- and why the biconditional it was framed as is false even over groups}

\author{MacBeth}
\date{2026-10-22}

\begin{document}
\maketitle

\begin{center}
\fbox{\begin{minipage}{0.92\textwidth}
\textbf{Recipient:} Rick \quad (CC: Robin)\\[2pt]
\textbf{Author:} MacBeth \qquad \textbf{Date:} 2026-10-22\\[2pt]
\textbf{Content commit:} \texttt{PENDING-STAMP}\\[4pt]
\small The decisive replacement for the vacuous Proposition~5.1 of the positive-home note
(\texttt{2026-10-19-handedness-positive-home}), requested by Rick's referee report (M3/Q2,
artifact \texttt{proofs/reviews/2026-10-21-triangulation-rick.md}). Registers in
\texttt{sesquilaw-magmoid-karoubi.json} under node
\texttt{positive-home-one-sided-absorbing-idempotent}, replacing child
\texttt{principal-normalizer-always-associative}.
\end{minipage}}
\end{center}

\bigskip

\section{What this proves, and the correction it carries}

Rick identified the structural ``why'' behind Theorem~B of the positive-home note ($\NK$
associative $\iff$ the normalizer $K$ is two-sided absorbing): it is a \textbf{closure
criterion}. This note states and proves it cleanly, in the $(\mathrm{LA})$-normalizer notation, as
an explicitly attributed specialization of the classical transversal / right-loop construction.

Two results, one positive and one corrective:

\begin{itemize}[leftmargin=1.4em]
\item \textbf{Theorem~1 (closure criterion --- the decisive statement).} If $\Im K$ is closed
  under composition $\cdot$, then $\# = \cdot|_{\Im K}$ and $\NK$ is associative. Contrapositive
  --- \emph{the decisive statement Rick asked for}: \textbf{if $\NK$ is handed
  (non-associative), then $\Im K$ is NOT closed under $\cdot$.} Airtight; it uses only
  idempotence of $K$ and associativity of $\Acat$.

\item \textbf{Theorem~2 (the converse FAILS --- sharpness, and a correction to the record).} The
  converse ``$\Im K$ not closed $\Rightarrow$ handed'', equivalently ``$\NK$ associative
  $\Rightarrow \Im K$ closed'', is \textbf{false} --- and it fails \emph{even over invertible
  (group) bases}. Hence the biconditional stated as the target in \texttt{state/PROVE.md}
  (``handedness $\iff (\mathrm{LA}) \wedge$ non-closure'') and the memory attribution
  ``$\#$ associative $\iff \Im K$ closed = classical right-loop group-iff-subgroup'' are
  \textbf{incorrect}: that biconditional is not a theorem. Non-closure is \textbf{necessary but not
  sufficient} for handedness. The honest content is the one-directional Theorem~1.
\end{itemize}

The genuine biconditional sitting underneath is \textbf{$\Im K$ closed $\iff \# = \cdot|_{\Im K}$}
(Prop~3) --- closure is equivalent to the normalized product \emph{coinciding} with the ambient
product, not to associativity.

\paragraph{Grades.} Theorem~1 and Prop~3 are \textbf{proved} (elementary, by hand, re-derived below
and verified on an exact finite model). Theorem~2's refutation is \textbf{proved} by two explicit
exact counterexamples (a minimal order-3 non-invertible monoid; and $C_4$/$D_4$/$Q_8$
normal-subgroup transversals), computed exactly in \texttt{fractions}/finite enumeration.
Attribution is to the classical transversal literature (Baer, Lal, Kiechle, Grabowski), with the
explicit note that the folklore \emph{biconditional} those sources are sometimes quoted for does
not hold.

\section{Setting and the standing definition}

The definitions are those of the positive-home note (\S1), repeated for self-containment.

\begin{definition}[normalizer, normalized magmoid]
Let $(\Acat, \cdot, 1)$ be a category (composition diagrammatic: $f\cdot g$ = ``$f$ then $g$'').
A \textbf{normalizer} is a family $K = K_{X,Y} : \Acat(X,Y) \to \Acat(X,Y)$ with
\begin{itemize}[leftmargin=2em]
\item \textbf{(K1)} $K$ idempotent: $K(K(a)) = K(a)$;
\item \textbf{(U)} $K(1_X) = 1_X$;
\item \textbf{(LA)} left-absorption: $K(K(x)\cdot c) = K(x\cdot c)$ for all composable $x, c$.
\end{itemize}
The \textbf{normalized magmoid} $\NK$ has $\Ob \NK = \Ob \Acat$,
$\NK(X,Y) := \Im K_{X,Y} = \{ a : K(a) = a \}$, product $f \# g := K(f\cdot g)$, units $1_X$.
$K$ is \textbf{two-sided} if $(\mathrm{RA}_\#)$: $K(f\cdot K(g\cdot h)) = K(f\cdot g\cdot h)$ for all
$f,g,h \in \Im K$. $\NK$ is \textbf{handed} if it is non-associative (equivalently, under
$(\mathrm{LA})$, $(\mathrm{LA}) \wedge \neg(\mathrm{RA}_\#)$ --- Theorem~B of the positive-home note).
\end{definition}

By (K1), $\Im K = \{a : K(a)=a\}$ (if $a = K(x)$ then $K(a)=K(K(x))=K(x)=a$; conversely immediate).
We say \textbf{$\Im K$ is closed under $\cdot$} if $f, g \in \Im K$ (composable) $\Rightarrow f\cdot
g \in \Im K$, i.e. $K(f\cdot g) = f\cdot g$.

\paragraph{Attribution (read before the theorem).} This whole package is the classical
\textbf{transversal / right-loop} construction. For a subgroup $H \le G$ with right transversal $T$
(one representative per right coset $Hg$, with $1 \in T$), the representative map $r : G \to T$,
$r(g) =$ the $t \in T$ with $Hg = Ht$, is a retraction onto $T$ satisfying exactly $(\mathrm{K1})$,
$(\mathrm{U})$, $(\mathrm{LA})$; the induced product $t \circ t' := r(t\cdot t')$ is our $\#$, and
$T = \Im r$. Sources: \textbf{Baer, ``Nets and groups,'' Trans.\ AMS 46 (1939) 110--141};
\textbf{Lal, ``Transversals in groups,'' J.\ Algebra 181 (1996) 70--81};
\textbf{Kiechle, \emph{Theory of K-Loops}, Springer LNM 1778 (2002)}; \textbf{arXiv:1309.7729, ``On
the Congruences in Right Loops''}; oidified: \textbf{Grabowski, ``An introduction to loopoids,''
arXiv:1510.03395} (Def~4.1 / Prop~4.3: $t\bullet t' = \pi(t\cdot t')$). The forward direction of
Theorem~1 (``$T$ a subgroup $\Rightarrow \circ = \cdot|_T$'') is classical and elementary. The
\emph{converse} --- the folklore biconditional ``right loop is a group $\iff$ transversal is a
subgroup'' --- is \textbf{not} a theorem (\S3).

\section{Theorem 1 --- the closure criterion}

\begin{theorem}[closure criterion]
Let $K$ be a normalizer on $\Acat$. If $\Im K$ is closed under $\cdot$, then:
\begin{enumerate}[label=(\arabic*),leftmargin=2.2em]
\item $f \# g = f\cdot g$ for all composable $f, g \in \Im K$ (i.e. $\# = \cdot|_{\Im K}$);
\item $\NK$ is associative.
\end{enumerate}
\noindent\textbf{Contrapositive (the decisive statement).} If $\NK$ is handed, then $\Im K$ is
\textbf{not} closed under $\cdot$.
\end{theorem}

\begin{proof}
(1) Let $f, g \in \Im K$ be composable. Closure gives $f\cdot g \in \Im K$, so $K(f\cdot g) =
f\cdot g$. Hence $f \# g = K(f\cdot g) = f\cdot g$.

(2) Let $f, g, h \in \Im K$ be composable. By (1), $f\#g = f\cdot g$; closure gives $f\cdot g \in
\Im K$, so with $h \in \Im K$, (1) applies again: $(f\#g)\#h = (f\cdot g)\#h = (f\cdot g)\cdot h$.
Symmetrically, $g\#h = g\cdot h \in \Im K$, so $f\#(g\#h) = f\#(g\cdot h) = f\cdot(g\cdot h)$. By
associativity of $\cdot$ in $\Acat$, $(f\cdot g)\cdot h = f\cdot(g\cdot h)$. Therefore
$(f\#g)\#h = f\#(g\#h)$.

The contrapositive is immediate.
\end{proof}

\begin{remark}[the proof uses almost nothing]
Theorem~1 invokes only \textbf{(K1)} (to read $\Im K$ as the fixed-point set) and
\textbf{associativity of $\cdot$ in $\Acat$}. It does \emph{not} use $(\mathrm{U})$ or even
$(\mathrm{LA})$. So closure is a sufficient condition for associativity that is insensitive to the
left-absorption law --- a first sign that it is strictly stronger than what associativity needs
(\S3). In the standing setting $(\mathrm{LA})$ holds, so ``handed'' means $(\mathrm{LA}) \wedge
\neg(\mathrm{RA}_\#)$; Theorem~1 gives $(\mathrm{LA}) \wedge \neg(\mathrm{RA}_\#) \Rightarrow$ not
closed. This is exactly Rick's ``handed $\Rightarrow \Im K$ not closed.''
\end{remark}

\begin{remark}[relation to Theorem~B]
Theorem~1 is the structural ``why'' behind Theorem~B of the positive-home note. Indeed closure
$\Rightarrow \NK$ associative $\iff (\mathrm{RA}_\#)$ (Theorem~B). So \textbf{closure $\Rightarrow
(\mathrm{RA}_\#)$}. The content of \S3 is that this implication is \emph{strict}:
$(\mathrm{RA}_\#)$ can hold without closure.
\end{remark}

\section{Theorem 2 --- the converse fails, even for groups}

The target in \texttt{state/PROVE.md} asked for the biconditional ``handedness $\iff$
non-closure'', equivalently ``$\NK$ associative $\iff \Im K$ closed''. Only the $\Leftarrow$ of the
second form (= Theorem~1) is true. The $\Rightarrow$ is false.

\begin{theorem}[sharpness]
There exist normalizers $K$ (satisfying (K1), (U), (LA)) with $\Im K$ \textbf{not} closed under
$\cdot$ and yet $\NK$ \textbf{associative}. Such examples exist already (a)~on a 3-element monoid
--- the minimal order --- with a non-invertible element, and (b)~on finite \textbf{groups} $C_4,
D_4, Q_8$. Hence non-closure is \textbf{necessary but not sufficient} for handedness, and
invertibility of the base does \textbf{not} restore the converse.
\end{theorem}

\subsection{A minimal non-invertible witness (order 3)}

Let $M = \{e, a, b\}$ be the monoid generated by one element $a$ with $a^3 = a$ (so $b := a^2$).
Its table (rows = left factor):
\[
\begin{array}{c|ccc}
\cdot & e & a & b \\\hline
e & e & a & b \\
a & a & b & a \\
b & b & a & b
\end{array}
\]
$M$ is associative with unit $e$ (both verified exactly). The element $a$ is \textbf{not
invertible} ($a\cdot x = e$ has no solution; likewise $x\cdot a = e$), so $M$ is \textbf{not a
group}.

Define the normalizer $K$ by $K(e) = e$, $K(a) = a$, $K(b) = a$. Then (checked over all $3\times3$
pairs, exactly):
\begin{itemize}[leftmargin=1.6em]
\item $(\mathrm{K1})$ $K\circ K = K$, $(\mathrm{U})$ $K(e)=e$, and $(\mathrm{LA})$ $K(K(x)\cdot c) =
  K(x\cdot c)$ hold \textbf{universally}.
\item $\Im K = \{e, a\}$, and it is \textbf{not closed}: $a\cdot a = b \notin \Im K$.
\item $\NK$ on $\{e,a\}$ has table $e\#e=e$, $e\#a=a\#e=a$, $a\#a=K(b)=a$ --- the 2-element
  semilattice, which is \textbf{associative} (all 8 triples).
\end{itemize}

So $(\mathrm{LA})$ holds, $\Im K$ is not closed, yet $\NK$ is associative: the converse of
Theorem~1 is refuted. An exhaustive search over all monoids of order $\le 3$ and all normalizers on
them shows \textbf{no such example exists at order $\le 2$} (0 pairs), and exactly \textbf{6} exist
at order 3 (across 3 monoid tables) --- so order 3 is minimal.

\paragraph{What goes wrong for the converse here.} $\Im K = \{e,a\}$ fails closure only at the
single product $a\cdot a = b$, whose normalization $K(b) = a$ ``lands back'' compatibly: the one
escaping product does not propagate to any 3-fold associativity defect. Non-closure is a
\emph{local} (binary) defect; handedness is a \emph{global} (ternary) defect, and the former does
not force the latter.

\subsection{Group witnesses ($C_4$, $D_4$, $Q_8$) --- invertibility does not help}

Take $K = r$, the right-coset representative map of a subgroup $H \le G$ with transversal $T$ (so
$\Im K = T$, $\# = \circ$). $(\mathrm{LA})$ holds in every case. The biconditional ``$\circ$
associative $\iff T$ subgroup'' \textbf{fails}:

\begin{itemize}[leftmargin=1.6em]
\item \textbf{$C_4 = \mathbb{Z}/4$, $H = \{0, 2\}$, $T = \{0, 1\}$.} Representatives:
  $r(0)=r(2)=0$, $r(1)=r(3)=1$. $\Im K = \{0,1\}$ is \textbf{not closed} ($1+1 = 2 \notin T$). But
  $(T,\circ)$: $1\circ 1 = r(2) = 0$, so $(T,\circ) \cong C_2$ is a \textbf{group} --- associative.
\item \textbf{$D_4$} (order 8): 12 such transversals; \textbf{$Q_8$}: 20, e.g. $H = \{\pm1\}$,
  $T = \{1,-i,-j,-k\}$ with $(-i)(-j) = k \notin T$, yet $(T,\circ) \cong C_2\times C_2$
  associative.
\end{itemize}

\paragraph{Mechanism (clean and general).} Whenever $H \triangleleft G$ is normal, the quotient map
restricts to a bijection $T \to G/H$, $t \mapsto Ht$, under which $\circ$ becomes the group law of
$G/H$: $r(t\cdot t') \mapsto H(t\cdot t') = (Ht)(Ht')$. So $(T, \circ) \cong G/H$ is a group for
\textbf{every} transversal of a normal subgroup --- while a transversal of a nontrivial normal
subgroup is almost never a subgroup. This is a whole family of counterexamples; the vanishing of
anomalies for $C_2, C_3, C_2\times C_2, S_3$ in the sweep is an accident of those small subgroup
lattices (every transversal of each normal subgroup there happens to be a subgroup), not a theorem.

This is the decisive correction: \textbf{invertibility does not rescue the converse.} The folklore
``right loop is a group $\iff$ transversal is a subgroup'' is simply false in the presence of a
nontrivial normal subgroup. The classical literature's \emph{true} content is the forward direction
(``subgroup $\Rightarrow \circ = \cdot|_T$'') --- i.e.\ Theorem~1 --- not the biconditional.

\section{Proposition 3 --- the biconditional that \emph{is} true}

\begin{proposition}
For a normalizer $K$: $\Im K$ is closed under $\cdot$ \textbf{if and only if} $\# = \cdot|_{\Im K}$
(i.e.\ $f\#g = f\cdot g$ for all composable $f,g \in \Im K$).
\end{proposition}

\begin{proof}
($\Rightarrow$) Theorem~1(1). ($\Leftarrow$) If $f\#g = f\cdot g$ for all $f,g \in \Im K$, then
$f\cdot g = f\#g = K(f\cdot g) \in \Im K$, so $\Im K$ is closed.
\end{proof}

So the honest ``iff'' pairs closure with \textbf{coincidence of the two products}, not with
associativity. Associativity is a strictly weaker ternary condition:
\[
\begin{array}{c}
\Im K \text{ closed} \ \iff\ \# = \cdot|_{\Im K} \qquad (\text{Prop 3}) \\[4pt]
\Downarrow \ \text{(strict, by Thm 2)} \\[4pt]
(\mathrm{RA}_\#) \text{ holds} \ \iff\ \NK \text{ associative} \qquad (\text{Theorem B})
\end{array}
\]
The classical transversal statement, stated correctly, is the top line specialized to groups:
``$T$ is a subgroup $\iff \circ$ is the restriction of $\cdot$.'' The middle-arrow strictness is
the whole content of \S3.

\begin{remark}[the obstruction, and the resonance with the weld $H^2$ program]
For $x, y \in \Im K$ write the de-normalization $x\cdot y = \varphi(x,y) \cdot (x\#y)$ where
$\varphi(x,y)$ is whatever $K$ strips off ($K(x\cdot y) = x\#y$). Then closure $\iff \varphi \equiv
1$ (trivial factor). Associativity of $\#$ is a \textbf{2-cocycle-type equation} on $\varphi$,
which $\varphi$ can satisfy \emph{without} being trivial --- precisely the normal-subgroup case,
where $\varphi$ is (the data of) a split extension and $(T,\circ) \cong G/H$. This is why closure
$\Rightarrow$ associativity but not conversely, and it is the same split-vs-nonsplit dichotomy that
the weld obstruction lives in (\texttt{h2-weld-class-is-splitting-not-associativity}): closure is
the \emph{trivial-factor} condition, strictly below the \emph{cocycle} condition. The handedness
home $\mathrm{Norm}$ of (\texttt{2026-10-19-handedness-positive-home}) realizes the strict gap from
the other side --- there $\Im K$ is not closed \textbf{and} $(\mathrm{RA}_\#)$ fails.
\end{remark}

\section{The corrected minimality statement (replacing Prop 5.1)}

Rick showed the old Proposition~5.1 (``a principal normalizer $K(x)=\kappa\cdot x$ by an idempotent
\emph{element} is always associative'') is vacuous under $(\mathrm{U})$: $K(1)=\kappa$ forces
$\kappa=1$, $K=\mathrm{id}$. The \textbf{decisive replacement} is:

\begin{corollary}[home of handedness]
If $\NK$ is handed then $\Im K$ is not closed under $\cdot$ (Theorem~1). Consequently the home of
handedness is a normalizer whose image is \textbf{not} a subcategory of $\Acat$. Non-closure is
\emph{necessary}; by Theorem~2 it is \emph{not sufficient}, so the full characterization remains
$(\mathrm{LA}) \wedge \neg(\mathrm{RA}_\#)$ (Theorem~B), of which non-closure is the clean necessary
shadow.
\end{corollary}

The $\mathrm{Norm}$ instance realizes it: $\Im K =$ row-stochastic kernels is \textbf{not closed}
under matrix product in the sense required --- the product of two normalized kernels de-normalizes
(mass is destroyed at $\bot$ states before re-normalization), so $K(f\cdot g) \neq f\cdot g$ in
general, and additionally $(\mathrm{RA}_\#)$ fails (exact witness, re-verified 2026-10-21:
$(f\#g)\#h = [8/11,3/11] \neq [2/3,1/3] = f\#(g\#h)$). Non-closure is the necessary half; the
sufficient defect is the ternary $(\mathrm{RA}_\#)$ failure.

\section{Verification (computational, exact)}

Scripts under \texttt{scratch/}, all exact (\texttt{fractions.Fraction} / finite enumeration),
2026-10-21:
\begin{itemize}[leftmargin=1.6em]
\item \texttt{2026-10-21-task1-order3-converse-counterexample.py}: the order-3 witness ---
  $(\mathrm{LA})$ universal, $\Im K$ not closed ($a\cdot a=b$), $\NK$ associative. Converse
  refuted.
\item \texttt{2026-10-21-task2-minimal-order-search.py}: exhaustive order $\le3$ --- 0 witnesses at
  order $\le2$, 6 at order 3. Minimality confirmed.
\item \texttt{2026-10-21-task3-group-converse-transversals.py}: $C_4/D_4/Q_8$ normal-subgroup
  transversals with $(\mathrm{LA})$ holding, $\circ$ associative, $T$ not a subgroup. Biconditional
  refuted over groups. (Mechanism $(T,\circ)\cong G/H$ confirmed.)
\item \texttt{2026-10-21-task4-norm-witness.py}: $\mathrm{Norm}$ handedness witness re-verified
  exactly.
\end{itemize}

\section{Gaps / honest scope}

\begin{itemize}[leftmargin=1.6em]
\item \textbf{Theorem~1 and Prop~3 are the proved core} (elementary, by hand; $(\mathrm{LA})$ not
  even required for Theorem~1). These are the registrable results.
\item \textbf{Theorem~2} is a refutation by explicit exact examples; the general ``normal $H
  \Rightarrow (T,\circ)\cong G/H$'' mechanism is classical (quotient group), stated as the reason,
  not a new theorem.
\item \textbf{The record correction} (the folklore ``group $\iff$ subgroup'' biconditional is
  false; the memory note and \texttt{PROVE.md} target over-stated it) is established by \S3 and
  should propagate to the registry and to
  \texttt{feedback-handedness-is-classical-transversal-not-original}.
\item No claim of originality beyond: (i)~the base is a \textbf{general associative category/monoid
  without invertibility} (the order-3 witness is non-invertible); (ii)~the clean \textbf{statement
  of the strict gap} closure $\subsetneq (\mathrm{RA}_\#)$ and its refutation of the folklore
  biconditional \emph{even over groups}; (iii)~the $\mathrm{Norm}$ substochastic-kernel instance.
  The closure criterion itself (forward direction) is classical, cited above.
\end{itemize}

\end{document}
